\documentclass[11pt,letterpaper]{report}
\usepackage[utf8]{inputenc}
\usepackage[T1]{fontenc}
\usepackage[spanish,es-tabla]{babel}
\usepackage[letterpaper,left=27mm,right=23mm,top=30mm,bottom=20mm,headheight=38pt,headsep=12pt]{geometry}
\usepackage[scaled]{helvet}
\renewcommand{\familydefault}{\sfdefault}
\usepackage{amsmath,amssymb,graphicx,booktabs,longtable,tabularx,array,enumitem,xcolor,fancyhdr,titlesec}
\usepackage{hyperref}
\usepackage{needspace}
\usepackage{etoolbox}
\pretocmd{\section}{\Needspace{5\baselineskip}}{}{}
\definecolor{azul}{RGB}{35,91,133}
\definecolor{rojo}{RGB}{162,39,50}
\definecolor{gris}{RGB}{70,82,92}
\hypersetup{colorlinks=true,linkcolor=azul,urlcolor=azul,pdfauthor={Aquiles Taylor Ramos Yapo},pdftitle={Proyecto académico de diseño de minicentral hidroeléctrica Limbani - R03}}
\urlstyle{same}
\setlength{\parindent}{0pt}
\setlength{\parskip}{3pt}
\setlength{\tabcolsep}{4pt}
\renewcommand{\arraystretch}{1.18}
\setlist{nosep,leftmargin=1.5em}
\setcounter{tocdepth}{1}
\setcounter{secnumdepth}{2}
\titleformat{\chapter}[display]{\normalfont\bfseries\color{azul}}{\large CAPÍTULO \thechapter}{6pt}{\LARGE}
\titlespacing*{\chapter}{0pt}{0pt}{14pt}
\titleformat{\section}{\large\bfseries\color{rojo}}{\thesection}{.7em}{}
\titleformat{\subsection}{\normalsize\bfseries\color{azul}}{\thesubsection}{.7em}{}
\fancyhf{}
\fancyhead[L]{\small\bfseries MINICENTRAL HIDROELÉCTRICA LIMBANI\\\footnotesize Centrales Hidroeléctricas | Primera unidad}
\fancyhead[R]{\footnotesize UNA PUNO\\R03 | 14.10.2026}
\fancyfoot[L]{\scriptsize Aquiles Taylor Ramos Yapo | 227368 | GRUPO A}
\fancyfoot[R]{\scriptsize Proyecto académico | \thepage}
\renewcommand{\headrulewidth}{.5pt}
\renewcommand{\footrulewidth}{.3pt}
\pagestyle{fancy}
\fancypagestyle{plain}{\fancyhf{}\fancyhead[L]{\small\bfseries MINICENTRAL HIDROELÉCTRICA LIMBANI\\\footnotesize Centrales Hidroeléctricas | Primera unidad}\fancyhead[R]{\footnotesize UNA PUNO\\R03 | 14.10.2026}\fancyfoot[L]{\scriptsize Aquiles Taylor Ramos Yapo | 227368 | GRUPO A}\fancyfoot[R]{\scriptsize Proyecto académico | \thepage}}
\newcommand{\estado}[1]{\par\noindent\colorbox{azul!8}{\parbox{\dimexpr\linewidth-2\fboxsep}{\small\textbf{Condiciones del modelo.} #1}}\par}
\newcommand{\dato}[1]{\csname dato#1\endcsname}
\newcommand{\fuente}[1]{{\small\color{gris}\textit{Fuente / trazabilidad: #1}}\par}
\newcolumntype{L}[1]{>{\raggedright\arraybackslash}p{#1}}
\emergencystretch=2em
\makeatletter\setlength{\@fptop}{0pt}\makeatother
\input{generado/valores}

\begin{document}
\hypersetup{pageanchor=false}
\begin{titlepage}\centering
{\Large\bfseries UNIVERSIDAD NACIONAL DEL ALTIPLANO\par}
\vspace{3mm}
{\normalsize FACULTAD DE INGENIERÍA MECÁNICA ELÉCTRICA,\\ELECTRÓNICA Y SISTEMAS\par}
{\normalsize ESCUELA PROFESIONAL DE INGENIERÍA MECÁNICA ELÉCTRICA\par}
\vspace{7mm}\includegraphics[height=27mm]{figuras/escudo-unap.png}\par
\vspace{7mm}{\large\bfseries CENTRALES HIDROELÉCTRICAS\par}
\textcolor{azul}{\rule{\textwidth}{1pt}}\par\vspace{5mm}
{\LARGE\bfseries\color{azul} PROYECTO ACADÉMICO DE DISEÑO\\DE UNA MINICENTRAL HIDROELÉCTRICA EN LIMBANI\par}
\vspace{5mm}
{\large Replanteamiento del aprovechamiento y suministro aislado\\a Limbani y Huancasayani\par}
\vspace{4mm}
{\normalsize Primera unidad: recurso hídrico, demanda, dimensionamiento,\\regulación, seguridad y evaluación de alternativas\par}
\vspace{8mm}
\begin{tabular}{L{4.1cm}L{9cm}}
\textbf{Estudiante} & Aquiles Taylor Ramos Yapo\\
\textbf{Código de estudiante} & 227368\\
\textbf{Docente} & Mg. Ing. Gregorio Meza Marocho\\
\textbf{Sección} & GRUPO A\\
\textbf{Ubicación} & Limbani, Sandia, Puno, Perú\\
\textbf{Revisión} & R03 - 14 de octubre de 2026\\
\textbf{Elaboración} & Trabajo académico estudiantil\\
\end{tabular}
\vfill\textcolor{azul}{\rule{\textwidth}{.5pt}}\par
{\small PUNO - PERÚ | 2026\par}
\end{titlepage}

\hypersetup{pageanchor=true}
\pagenumbering{roman}
\chapter*{Resumen}
\addcontentsline{toc}{chapter}{Resumen}
Se desarrolla un \textbf{proyecto académico de diseño de una minicentral hidroeléctrica, elaborado por un estudiante}, para el curso de Centrales Hidroeléctricas. Se replantea el aprovechamiento existente referido en Limbani, Sandia, Puno, y se considera suministro \textbf{aislado} a Limbani y Huancasayani. Las referencias de 200 kW, salto cercano a 50 m y turbina Michell-Banki constituyen antecedentes por verificar; no se atribuyen a una placa inspeccionada.

La investigación recuperó series públicas PISCO de SENAMHI: \textbf{16\,436 días y 540 meses entre 1981 y 2025}. Se identificó el tramo COMID 9080887 y se calculó una cuenca exploratoria en la captación propuesta. El ajuste por área utiliza superficie geodésica, corrigiendo la distorsión de Web Mercator del atributo original. La serie transferida es una \textbf{estimación sin calibración local}; la estación hidrológica pública de Limbani reporta niveles y no se recuperó su curva nivel--caudal. El contraste entre productos y la delimitación digital muestran incertidumbres que se conservan en el diseño.

El censo INEI 2017 distingue cuatro registros del ámbito adoptado: Limbani, Huancasayani Centro, Bellavista y Lince. Suman \dato{PopBase}\ habitantes y \dato{VivsBase}\ viviendas ocupadas censales. La proyección de estudio y las calificaciones de la norma rural de 2023 producen \dato{Users}\ conexiones potenciales y una envolvente de demanda de \dato{Demand}\ kW. Estos valores no son un padrón ni una demanda medida. Por esa cobertura conservadora, se comparan 150, 200, 400 y 600 kW, y se dimensiona una alternativa modular de \textbf{tres grupos de 200 kW}, con respaldo aislado. La ampliación no se impone como inversión óptima: depende del padrón, el recurso y los costos.

La propuesta cartográfica usa salto bruto de estudio \dato{Hb}\ m, tubería de \dato{LongPipe}\ m y diámetro interior de \dato{Di}\ mm. Con caudal nominal calculado de \dato{Qnom}\ m$^3$/s entrega 600 kW en bornes y aproximadamente \dato{Pusers}\ kW en usuarios. Sin embargo, el escenario diario proporciona solo \dato{Pfirmstudy}\ kW con excedencia del 95\% y presenta generación nula en algunos periodos. La hidroeléctrica por sí sola \textbf{no demuestra servicio continuo de 150--200 kW durante el estiaje}.

Para el perfil de demanda supuesto, la energía hidroeléctrica servida se estima en \dato{EHidro}\ MWh/año, frente a \dato{Eload}\ MWh/año demandados. Se estudia respaldo modular, regulación de frecuencia/tensión, protección de red débil, transitorios y costos. El presupuesto conservador de reposición y ampliación alcanza aproximadamente S/\dato{CapexM}\ millones, con incertidumbre elevada y rentabilidad financiera desfavorable en los ingresos examinados. Se recomienda validar usuarios, reutilización de activos y caudales antes de escoger una inversión.

\estado{El alcance es académico. No se emite un expediente técnico definitivo ni una autorización para construir. Las comprobaciones posteriores se identifican por capítulo y en la matriz normativa. Se conservan las fuentes originales y los cálculos reproducibles.}
\begingroup\small\setlength{\parskip}{0pt}\tableofcontents\endgroup
\clearpage\pagenumbering{arabic}
\chapter{Introducción, objetivos y alcance académico}
\section{Introducción}
Diseñar una minicentral exige relacionar el recurso hídrico con la población atendida, el salto aprovechable y la capacidad real de las máquinas. En un sistema aislado, la frecuencia y la tensión dependen de sus propios generadores; el déficit de agua no puede compensarse importando del SEIN. El diseño debe incorporar regulación, respaldo o gestión de carga y explicar cómo cambia el servicio durante el año.

La instrucción del docente fue diseñar una minicentral hidroeléctrica, sin una rúbrica técnica detallada proporcionada. Este trabajo define un alcance académico integrado: investigación pública, modelación, selección argumentada y comprobaciones de ingeniería con sus límites. Las normas aplicables orientan las verificaciones; el estudio no reemplaza los permisos, los ensayos ni la responsabilidad profesional de una obra real.

El proyecto parte del borrador recibido y del trabajo previo de redes de distribución de Aquiles. R03 modifica el modo de operación que R02 había supuesto conectado a red. Retira el objetivo rígido de exportación de 150 kW, introduce población, caudales públicos simulados y una comparación de tamaños. El rango 150--200 kW se conserva como alternativa inicial, abierto a cambios justificados.

\section{Objetivo general}
Desarrollar un proyecto académico de diseño para replantear el aprovechamiento hidroeléctrico de Limbani y atender, mediante un sistema aislado, la demanda de Limbani y Huancasayani, identificando las restricciones del recurso, las condiciones de seguridad y la viabilidad de las alternativas.

\section{Objetivos específicos}
\begin{enumerate}
\item Identificar fuentes ANA, SENAMHI, INEI y antecedentes cartográficos; distinguir observación, simulación, estimación y decisión de diseño.
\item Analizar variación estacional, estiaje, reservas de agua, duración de caudales y disponibilidad de potencia/energía.
\item Construir una demanda por localidad y contrastarla con tamaños de planta, sin duplicar el total distrital ni usar huellas sintéticas como usuarios reales.
\item Dimensionar obras hidráulicas, grupo, tubería, auxiliares y sistema eléctrico; comprobar pérdidas y transitorios con hipótesis trazables.
\item Diseñar la regulación y las funciones de protección para operación aislada y describir previsiones para una eventual interconexión futura.
\item Elaborar metrados, costos, planos académicos y conclusiones coherentes con los cálculos y la matriz normativa.
\end{enumerate}

\section{Ficha de alcance}
\begin{longtable}{L{4.2cm}L{10cm}}
\toprule Elemento & Decisión y condición\\\midrule
Naturaleza & Proyecto académico estudiantil de diseño; replanteamiento de un aprovechamiento existente referido por el estudiante\\
Beneficiarios & Limbani y agrupación Huancasayani en el distrito de Limbani; Centro, Bellavista y Lince se distinguen por código censal\\
Operación inicial & Sistema aislado, 60 Hz; sin enlace ni importación del SEIN\\
Potencia & Comparación 150, 200, 400 y 600 kW; alternativa de cobertura conservadora 3 por 200 kW, pendiente de padrón y aforos\\
Obras & Derivación de pasada, desarenado, canal, cámara, tubería, casa de máquinas y restitución\\
Electricidad & 380 V de generación, transformación a 10 kV y distribución 380/220 V como propuesta académica coherente con el antecedente\\
Continuidad & Hidro más respaldo y prioridad de cargas; no se atribuye suministro permanente a la placa hidroeléctrica\\
Horizonte & Diseño de escenario 2026, sensibilidad de crecimiento y previsión de eventual SEIN entre 2046 y 2076\\
Entregables & Memoria y anexos, planos/esquemas académicos, tablas reproducibles y material de exposición\\
\bottomrule\end{longtable}

\section{Método y trazabilidad}
El directorio de datos guarda entradas con unidades, procedencia y estado. Los originales recibidos y los archivos descargados se conservan inalterados. Los resultados se generan a partir de esas entradas; las tablas no se ajustan manualmente para producir una potencia deseada. Las cantidades calculadas alimentan el presupuesto y los esquemas.

\begingroup\small
\begin{longtable}{L{3.2cm}L{6cm}L{4.8cm}}
\toprule Tipo de información & Ejemplo del trabajo & Alcance\\\midrule
Dato documental & Población INEI 2017, coordenadas de una ficha ANA & Verificado en su fuente; no se convierte en dato actual o de otra instalación\\
Simulación pública & Caudales PISCO ARNOVIC y GR2M & Producto modelado; no registro de aforos\\
Estimación del proyecto & Transposición por área, demanda proyectada, costos & Método y supuestos expresos; calibración pendiente\\
Decisión académica & Configuración aislada, diámetros, equipos y regulación & Propuesta para comparar y comprobar\\
Dato no disponible & Padrón activo, resistividad, geotecnia y placa existente & Se registra como pendiente; no se sustituye por una medición ficticia\\
\bottomrule\end{longtable}\endgroup

\chapter{Ubicación, antecedentes y emplazamiento}
\section{Ubicación política y referencias territoriales}
El ámbito pertenece al distrito de Limbani, provincia de Sandia, departamento de Puno. La unidad hidrográfica identificada en resoluciones ANA locales es Inambari, código 46648 [H07, H10]. La propuesta utiliza la vertiente amazónica; no se transfieren automáticamente caudales de ríos del Titicaca por encontrarse también en Puno.

\begin{figure}[htbp]\centering
\begin{minipage}{.48\linewidth}\centering\includegraphics[width=\linewidth,height=6cm,keepaspectratio]{figuras/ubicacion-puno.jpg}\end{minipage}\hfill
\begin{minipage}{.48\linewidth}\centering\includegraphics[width=\linewidth,height=6cm,keepaspectratio]{figuras/ubicacion-sandia.jpg}\end{minipage}
\caption{Ubicación político-geográfica de Puno y Sandia. Imágenes conservadas del antecedente de Aquiles; fecha y autor originales no acreditados. Referencia territorial, sin escala del proyecto.}
\end{figure}

\section{Central existente y alcance del replanteamiento}
Se tomó como referencia una central de aproximadamente 200 kW, salto cercano a 50 m y turbina Michell-Banki, conforme a la indicación del estudiante y al proyecto de redes. No se recuperó una ficha oficial o fotografía legible de placa que confirme estos tres valores. Los 50 m no permiten distinguir salto bruto y neto; por ello se conserva esa incertidumbre y se construye una hipótesis independiente de geometría para los cálculos.

El trabajo de línea primaria y redes usa 200 kW como \emph{entrada}, no como resultado del diseño hidroeléctrico. Su cuestionario confirma el interés por 10 kV y control aislado. Se aprovechan esos antecedentes, pero no se importan como hechos sus afirmaciones de precisión catastral, demanda o caudal. El JSON menciona 764 viviendas y 2.15 km; la memoria y los planos usan otros conteos, incluidos 587 acometidas y una longitud cercana a 1.97 km. Esos datos son antecedentes del casco urbano y no cubren automáticamente Huancasayani.

La carpeta Catastro\_Limbani contiene 482 huellas en el avance original y una base complementada mediante algoritmos. Las huellas añadidas representan geometría propuesta; no son evidencia de ocupación o contratos eléctricos. La demanda R03 utiliza los registros censales, no la suma de polígonos de casas.

\section{Ficha de coordenadas y procedencia}
La base recibida declara WGS84 / UTM zona 19 Sur, \textbf{EPSG:32719}, en metros. Se transforma a EPSG:4326 solo para cruzar estaciones y productos públicos. Las alturas de mapas tienen referencia vertical por confirmar y no se presentan como nivelación de campo.

\begin{longtable}{L{4.2cm}L{4.4cm}L{5.5cm}}
\toprule Referencia & Coordenadas & Procedencia y certeza\\\midrule
Casa candidata CM-A & E 425580; N 8436580 m & Memoria del trabajo de redes; punto asignado académicamente, sin inspección\\
Casa candidata CM-B & E 425715; N 8436650 m & Script/plano del mismo antecedente; difiere de CM-A unos 152 m\\
Casa en borrador Gemini & E 425260; N 8441710 m & Aproximadamente 5 km al norte de los antecedentes urbanos; retirada del emplazamiento activo\\
Captación de estudio & E 425542; N 8436231 m, aproximadamente & Punto propuesto 500 m aguas arriba por recorrido del eje PISCO desde CM-B; no identifica la bocatoma existente\\
Estación HLG Limbani & Lat. -14.1481; lon. -69.6902 & Mapa SENAMHI, código 260501; no es casa de máquinas ni captación\\
Cámara y restitución & Sin coordenadas confirmadas & Se desarrollan por perfil y esquema académico; el trazo definitivo requiere levantamiento\\
Huancasayani Centro & Lat. -14.1917532; lon. -69.7048438 & OSM referencial, contrastado por nombre/distrito; no coordenada oficial INEI\\
Huancasayani Bellavista & Lat. -14.1983659; lon. -69.7159784 & OSM referencial; Lince no se ubicó con certeza en esta investigación\\
\bottomrule\end{longtable}

\begin{figure}[htbp]\centering\includegraphics[width=.76\linewidth,height=12cm,keepaspectratio]{figuras/emplazamiento-r03.png}
\caption{Base recibida y puntos de distinta procedencia. La escala gráfica permite leer distancias cartográficas; no acredita precisión del levantamiento ni ubicación de instalaciones existentes.}
\end{figure}

\section{Aprovechamiento topográfico y corredor de distribución}
El perfil SRTM90 consultado sobre el eje fluvial tiene elevaciones de 3285 m a 500 m aguas arriba y 3239 m en la referencia próxima a CM-B [H05]. Son elevaciones de celdas, no niveles de agua ni cotas del eje de turbina. Para el diseño académico se reserva 0.30 m para 150 m de canal a pendiente 0.002 y 0.40 m para transiciones y diferencias de nivel. Se plantea cámara a 3284.3 m y referencia de turbina a 3239 m, produciendo \dato{Hb}\ m de salto bruto de estudio. La cámara debe situarse en una ladera compatible con ese nivel; no se asegura tal ubicación a partir de un eje de río.

El recorrido fluvial de 500 m motiva una propuesta de canal de 150 m y tubería de 350 m. Esta longitud de tubería es una aproximación de recorrido tridimensional, que se estudia además entre 300 y 450 m. No se conserva la tubería de 175 m del borrador ni se aumenta el salto para obtener la potencia objetivo. La cartografía de 90 m no resuelve por sí sola el corredor de conducción, movimientos de tierra o estabilidad de ladera.

\begin{figure}[htbp]\centering\includegraphics[width=.9\linewidth,height=10cm,keepaspectratio]{figuras/corredor-r03.png}
\caption{Relación territorial de las localidades. Las líneas punteadas son vínculos de ubicación, no trazos de postes. El recorrido de 8 km usado en electricidad es una hipótesis de corredor, superior al enlace urbano del antecedente.}
\end{figure}

\estado{La casa de máquinas existente, la captación y la cámara no están identificadas con certeza. Los puntos de pueblo, estación y edificios candidatos mantienen etiquetas diferentes. La definición ejecutiva exige visita, placas, GPS, nivelación, servidumbres y faja marginal.}

\chapter{Recurso hídrico, estiaje y energía}
\section{Datos públicos recuperados y su naturaleza}
La estación hidrológica Limbani del mapa público del SENAMHI tiene código 260501, categoría HLG y condición de transmisión diferida. El visor permite identificar la estación, pero la descarga de datos está protegida por una verificación de usuario. No se obtuvo una serie de caudales ni su curva nivel--caudal. El boletín de diciembre de 2025 consultado informa \textbf{nivel en metros}; convertirlo en caudal sin una curva de gasto sería incorrecto. Tampoco se debe confundir esta estación con la meteorológica CO 114049. Sus coordenadas publicadas y los inventarios históricos no coinciden exactamente y se registran por separado (H03 y H17).

Se recuperaron íntegramente los productos públicos PISCO de SENAMHI, publicados en HydroShare: ARNOVIC v1.1 diario y GR2M v2.0 mensual (H01--H02). Ambos abarcan 1981--2025. Son \textbf{caudales simulados de modelos hidrológicos}, no aforos en Limbani. La extracción del COMID 9080887 contiene 16\,436 días consecutivos y 540 meses, en \(\mathrm{m^3/s}\), sin valores faltantes ni negativos. Los archivos NetCDF, sus vectores y sus huellas SHA-256 se conservan para reproducir el cálculo. La licencia publicada es CC BY 4.0; los autores y la procedencia se mantienen en el registro de fuentes.

El tramo cartográfico pasa a unos 25 m de la referencia CM-B, pero su caudal representa la \textbf{salida de una cuenca aguas abajo}. La casa de máquinas está aproximadamente al 40\% del recorrido del tramo de 32.6 km. Aplicar directamente ese caudal a la captación sobreestimaría el recurso. Por ello se estudia una transposición espacial explícita, con incertidumbre.

\section{Nombre de los ríos y antecedentes próximos}
La denominación río Limbani está respaldada por SENAMHI y los antecedentes locales. No se ha verificado que ``río Para'' sea el nombre oficial de un río aprovechable en el emplazamiento. Se investigó Phara, distrito vecino, sin declarar que Para y Phara sean equivalentes. La resolución ANA 370-2017 identifica una captación en el río \textbf{Cullucachi}, y otro antecedente ANA menciona el río Comunidad en Phara (H07 y H18). La licencia de la minicentral de Phara concede 90 l/s y 2\,838\,240 \(\mathrm{m^3/a}\) para 100 kW; se trata de un derecho de uso de \emph{otra} central, no de un registro natural de caudales ni de una autorización para Limbani.

La tesis de Goyzueta Arce (2017), de la Universidad Nacional del Altiplano, aporta un diseño de Phara con turbina Pelton, \(Q=0.24\ \mathrm{m^3/s}\) y salto neto de 134.25 m. Su cuadro de cinco años no identifica los años calendario. No se adopta como serie temporal oficial ni se transfiere mediante regresión: faltan aforos concurrentes, áreas verificadas, precipitación comparable y control de usos del agua. Su mayor salto tampoco puede trasladarse a Limbani.

\section{Crecidas y seguridad: análisis independiente}
El máximo diario simulado, \dato{Qhigh} \(\mathrm{m^3/s}\), no es pico instantáneo ni avenida de período de retorno conocido. Para bocatoma y casa se requieren extremos, secciones, niveles, socavación y sedimentos, con períodos de retorno justificados por riesgo. Este estudio de seguridad se desarrolla separado de generación; el máximo simulado no certifica las obras.

\section{Transposición por área y comprobaciones cartográficas}
Se delimitó una cuenca académica para la captación propuesta mediante un modelo de elevación público Mapzen, con relleno de depresiones y dirección de drenaje D8. El algoritmo conserva la acumulación de celdas y se verificó con pruebas independientes. La posición se ajusta a la celda de mayor acumulación dentro de 250 m: esta tolerancia reconoce que la captación todavía no está levantada en campo.

El área del atributo del vector PISCO es \(\dato{Amerc}\ \mathrm{km^2}\), pero coincide con el área calculada en \textbf{Web Mercator}. Esa proyección distorsiona las superficies. Se reemplaza por el área geodésica WGS84, \(A_r=\dato{Aref}\ \mathrm{km^2}\). La cuenca estimada en captación tiene \(A_c=\dato{Acapt}\ \mathrm{km^2}\). El factor de estudio resulta:
\begin{equation}
 Q_c(t)=\alpha Q_r(t),\qquad \alpha=\frac{A_c}{A_r}=\dato{Alpha}.
\end{equation}
Esta relación supone rendimientos específicos uniformes entre una subcuenca y su cuenca de referencia. No es una regresión ajustada ni una regionalización calibrada. Cerca del \(\dato{NestedPct}\%\) de la subcuenca delimitada está dentro del polígono PISCO, y sus áreas cambian menos del 1\% al repetir la delimitación a resoluciones de 30 y 90 m. Esto verifica estabilidad geométrica local, \textbf{no exactitud del caudal}.

Existen dos limitaciones relevantes. Primero, la cota media de la subcuenca es aproximadamente \dato{ElevationLocal} m, frente a \dato{ElevationRef} m en el polígono de referencia; la precipitación, evapotranspiración y almacenamiento pueden diferir. Segundo, el ajuste del DEM en la salida PISCO captura una confluencia ambigua y no reproduce satisfactoriamente el área de referencia. Ese control \textbf{no fue aceptado}. No se normaliza el caudal por una cuenca de salida D8 arbitraria. Se mantiene el área geodésica del producto y se examinan sesgos de 0.5, 0.7, 1.0 y 1.3; no son intervalos de confianza estadísticos.
\begin{figure}[htbp]\centering\includegraphics[width=.95\linewidth]{figuras/cuenca-r03.png}\caption{Cuenca propuesta y referencia PISCO. Delimitación de gabinete; no es una cuenca aforada ni un levantamiento de campo.}\end{figure}

\section{Estacionalidad y curva de duración}
\input{generado/hidro_estadisticas}
El caudal diario medio transpuesto es \dato{Qmean} \(\mathrm{m^3/s}\), y su Q95 diario es \dato{Qlow} \(\mathrm{m^3/s}\). Q95 significa excedencia aproximada durante el 95\% de los días simulados del período; no garantiza ese caudal cada año ni autoriza una extracción. La curva se construye ordenando los valores de mayor a menor y asignando posición de Weibull \(p=m/(n+1)\). Las cuantificaciones del cuadro se interpolan sobre la distribución empírica y conservan resolución diaria.
\begin{figure}[htbp]\centering\includegraphics[width=.94\linewidth]{figuras/caudales-r03.png}\caption{Serie transpuesta y variación mensual. Las dos líneas proceden de modelos públicos, no de mediciones en la captación.}\end{figure}
\begin{figure}[htbp]\centering\includegraphics[width=.90\linewidth]{figuras/duracion-r03.png}\caption{Duración diaria del caudal natural estimado y del caudal utilizable bajo reservas de estudio.}\end{figure}
La media mensual de GR2M difiere en \dato{MonthlyBias}\% de la media diaria agregada a meses de ARNOVIC; la correlación mensual es aproximadamente 0.76. Los modelos comparten información hidrometeorológica y no constituyen validación independiente. La suavización mensual oculta los días críticos. Su potencia con excedencia mensual de 95\%, \dato{PmonthlyLow} kW en usuarios, no reemplaza al indicador diario ni demuestra continuidad del servicio.

\section{Reserva ecológica, otros usos y balance de agua}
La RJ 267-2019-ANA distingue métodos y alcances para definir el caudal ecológico según la categoría ambiental y la coordinación con la AAA. No establece un porcentaje universal del 15\%. Para este ejercicio se reserva el 15\% de la media climatológica de cada mes, \textbf{como supuesto de sensibilidad}, hasta disponer del estudio hidroecológico y de la decisión de ANA. Se deja además una reserva ilustrativa de 0.005 \(\mathrm{m^3/s}\) para otros usos; tampoco es un inventario de derechos.
\begin{align}
 Q_e(t)&=\min[Q_c(t),\,0.15\overline{Q}_{c,mes}],\\
 Q_o(t)&=\min[0.005,\,\max(0,Q_c-Q_e)],\\
 Q_a(t)&=\max(0,Q_c-Q_e-Q_o),\\
 Q_c(t)&=Q_e(t)+Q_o(t)+Q_{t,\,posible}(t)+Q_{remanente}(t).
\end{align}
El caudal posible de turbinar se limita por equipos y caudal máximo de diseño. La serie exportada describe \textbf{capacidad potencial}, no una consigna real de extracción horaria. En despacho parcial el caudal realmente turbinado sería menor, y el excedente debe quedar en el río o descargarse por dispositivos dimensionados. La reserva ecológica se mantiene en el tramo derivado, sin sustituirla por la restitución aguas abajo.

La resolución ANA de 2019 para uso poblacional en Limbani reconoce 3.7 l/s de cuatro manantiales. No corresponde restar automáticamente ese volumen del río en esta captación: deben identificarse conexiones hidráulicas, titularidad y derechos concurrentes mediante el RADA y el análisis de disponibilidad hídrica. No se acreditan disponibilidad ni licencia obtenidas.

\section{Potencia y energía: comparación de alternativas}
Se utiliza \(P=\rho gQH_n\eta_t\eta_m\eta_g\), con pérdidas de conducción, eficiencias dependientes de carga, auxiliares y pérdidas eléctricas calculadas por separado. Las curvas parciales adoptadas son hipótesis académicas; se requieren curvas del proveedor para seleccionar equipos. La disponibilidad mecánica de 95\% se aplica a la \textbf{energía esperada}, sin convertirla en una garantía de potencia firme. La demanda se mantiene con el perfil 2026 para someterla a los 45 años climáticos simulados. Cada año se normaliza a 365 días.
\input{generado/alternativas}
\input{generado/hidro_meses}
El caso de 600 kW permite aproximadamente \dato{Pusers} kW en usuarios en régimen nominal, frente a \dato{Demand} kW de demanda máxima estimada. Sin embargo, el indicador diario de excedencia de 95\% es solamente \dato{Pfirmstudy} kW y el mínimo del modelo es \dato{PMin} kW. Hay \dato{DaysZero} días equivalentes por año sin generación bajo las reservas adoptadas. \textbf{Ni 150 ni 200 kW pueden ofrecerse como potencia continua en estiaje}. En julio, la media de potencia disponible es \dato{PJulMean} kW y el indicador diario de excedencia de 95\% es \dato{PJulLow} kW.

La capacidad hidroeléctrica anual disponible es \dato{Epotential} MWh, pero la energía que coincide con las cargas es \dato{EHidro} MWh. No se vende ni valora como ingreso la energía sobrante. El déficit hidroeléctrico es \dato{Edef} MWh/año; requiere respaldo o reducción de servicio. Incrementar de 400 a 600 kW cubre mejor el pico, aunque su ganancia de energía útil es pequeña frente al costo del tercer grupo. El escenario de 600 kW es una \textbf{envolvente conservadora de cobertura}, no la demostración de un óptimo económico.
\input{generado/sesgo}
\begin{figure}[htbp]\centering\includegraphics[width=.94\linewidth]{figuras/oferta-demanda-r03.png}\caption{Demanda sintética, capacidad variable y déficit de energía. No es un registro medido de operación.}\end{figure}

\chapter{Población atendida y demanda del sistema aislado}
\section{Ámbito censal y usuarios potenciales}
El Directorio Nacional de Centros Poblados de INEI, Censos 2017, tomo 4, páginas impresas 1713--1714, distingue cuatro registros relevantes: Limbani, Huancasayani Centro, Huancasayani Bellavista y Huancasayani Lince (H06). Son centros poblados en el directorio; no se inventa una categoría jurídica adicional. Se interpreta Huancasayani como el conjunto de esos tres registros para un escenario conservador, hasta que el padrón confirme el ámbito efectivo del servicio.

El distrito de Limbani registra 2\,970 habitantes y 2\,247 viviendas ocupadas; \textbf{ese total no se suma} a los centros poblados. Las cuatro localidades seleccionadas suman \dato{PopBase} habitantes y \dato{VivsBase} viviendas ocupadas en 2017. La nota censal incluye viviendas con personas ausentes o de uso ocasional. Por ello, ocupada no significa necesariamente residencia permanente, conexión activa o consumidor simultáneo. No se han contado 764 usuarios a partir del catastro generado ni se han duplicado viviendas.
\input{generado/poblacion}
Se verificó la existencia de resultados nacionales del Censo 2025, publicados en 2026, pero no se recuperó en esta revisión un directorio desagregado de estas localidades que sustituya la tabla anterior. El año base local utilizado sigue siendo 2017, no una población pretendidamente actual. La proyección geométrica adopta \(r=1\%\) anual y redondea cada localidad hacia arriba:
\begin{equation}N_{2026}=\lceil N_{2017}(1.01)^9\rceil.\end{equation}
Resultan \dato{PopProj} habitantes y \dato{Users} conexiones \emph{potenciales}, pendientes de padrón. Las tasas 0 y 2\% producirían límites aproximados de 946 y 1\,131 viviendas para el conjunto, antes del redondeo por localidad; esta sensibilidad no es una proyección oficial de INEI. La migración y ocupación estacional podrían reducir la demanda efectiva.

\section{Calificación eléctrica y simultaneidad}
La RD 108-2023-MINEM/DGE derogó la RD 031-2003 y aprobó nuevas bases para redes secundarias rurales con conductores autoportantes (N21). Sus numerales 3.2 y 4.7 utilizan 600 W/lote para localidad tipo I, 400 W/lote para tipo II y factor de simultaneidad doméstico 0.80. No se mantiene la antigua combinación de 350 W con simultaneidad 0.52.

Limbani se estudia como tipo I por su condición de capital distrital. Centro y Bellavista se asignan provisionalmente a tipo I, y Lince a tipo II; las tres clasificaciones necesitan inspección del trazado urbano y revisión con la entidad competente. Es posible adoptar calificaciones diferentes cuando estén sustentadas, pero no rebajarlas para hacer coincidir la demanda con 200 kW. El resultado doméstico es:
\begin{equation} P_{dom}=0.80\sum_i N_i p_i=\dato{Dom}\ \mathrm{kW}.\end{equation}
El alumbrado incluye 226 puntos LED de 50 W y pérdidas de 4.10 W por punto como referencia del cuadro de N21; resulta \dato{AP} kW con simultaneidad 1. Las cantidades son una previsión de diseño y deben contrastarse con inventario, disposición vial y fotometría. La carga especial se reserva con el siguiente inventario propuesto:
\input{generado/servicios}
\begin{equation}P_{max}=P_{dom}+P_{AP}+P_{especial}=\dato{Demand}\ \mathrm{kW}.\end{equation}
No se vuelve a aplicar simultaneidad a ese total. Se conserva el factor de potencia doméstico 1.0 de la base normativa y factores específicos de alumbrado y servicios. Los auxiliares de generación pertenecen a la central, por lo que se descuentan en el balance de oferta y no se duplican en la demanda de usuarios.

\section{Perfil, prioridades y crecimiento}
El perfil de 24 horas es sintético: representa carga nocturna baja, servicios diurnos y pico vespertino. No procede de mediciones de medidores. Su energía diaria es \dato{DailyE} kWh, la anual \dato{Eload} MWh y el factor de carga \dato{LoadFactor}\%. Se entrega el CSV horario para sustituirlo por lecturas reales. Los arranques de motores, desequilibrios y cargas no lineales requieren un estudio adicional; el perfil representa potencia activa estable.

Se propone un circuito prioritario de 131.6 kW en la hora de máxima carga, para salud, comunicaciones, agua y un suministro doméstico básico. Su continuidad exige respaldo: incluso esa reserva supera ampliamente el indicador hídrico diario de estiaje. La secuencia de reducción de carga debe preservar primero seguridad y salud, luego servicios básicos, y finalmente restringir actividades diferibles mediante acuerdos de operación y protección selectiva.
\input{generado/proyeccion}
El crecimiento posterior supone 1\% anual de usuarios y 1.5\% anual de carga por usuario, sin presentar esos porcentajes como pronóstico oficial. La capacidad inicial de 600 kW no cubre automáticamente esos picos futuros. Deben revisarse demanda real, eficiencia energética, incorporación por etapas y oferta híbrida antes de ampliar instalaciones. La posible llegada del SEIN entre 2046 y 2076 es una consideración de adaptación, no un supuesto que cierre el déficit actual.

\section{Resultado y selección académica}
Con el ámbito y las calificaciones conservadoras adoptadas, 150--200 kW son insuficientes para el pico. El caso de 400 kW permite mayor aprovechamiento que 200 kW y puede ser una etapa si el padrón demuestra una demanda menor o se acepta limitar cargas. Para estudiar cobertura completa del escenario se desarrolla 600 kW hidroeléctricos en tres grupos y respaldo diésel de 600 kW en tres módulos. La selección queda condicionada a disponibilidad hídrica, demanda real y evaluación económica. No se propone adquirir equipos sobre la base de una proyección de viviendas sin verificar.

\chapter{Disposición y dimensionamiento hidráulico}
\section{Alternativa de gabinete y perfil}
Se estudia una derivación de aproximadamente 500 m de corredor: captación y reja, canal de 150 m, desarenador, cámara de carga, tubería de presión de 350 m, casa de máquinas y restitución. La cota de cámara 3284.3 m se estima restando 0.3 m de caída del canal y 0.4 m de transiciones a una referencia de terreno de 3285 m. La referencia de casa de máquinas es 3239 m; resulta \(H_b=\dato{Hb}\) m. Estas cifras proceden de gabinete, no de nivelación de lámina de agua o eje de turbina.
\begin{figure}[htbp]\centering\includegraphics[width=.94\linewidth]{figuras/perfil-hidraulico-r03.png}\caption{Perfil funcional académico. La existencia de una ladera compatible con la cámara y el canal debe comprobarse; los segmentos no constituyen una rasante replanteada.}\end{figure}
Un trazado más largo de 450 m aumentaría la fricción frente a 350 m aproximadamente en proporción a su longitud para igual diámetro y caudal. Reducirla a 300 m puede disminuir pérdidas, pero exige un corredor físicamente disponible. Una variación de \(\pm5\) m en salto bruto implica aproximadamente \(\pm11\%\) de potencia hidráulica a igual caudal, antes de recalcular rendimiento y pérdidas. No se adopta el salto de 50 m como un valor medido ni se cambian cotas para obtener una potencia deseada.

\section{Captación y control de sólidos}
El caudal máximo hidráulico es \(Q_d=\dato{Qdesign}\ \mathrm{m^3/s}\). Se considera captación lateral con compuerta, reja inclinada, canal de limpieza y paso permanente de reserva ecológica. La toma lateral facilita separar la derivación del paso de crecidas. Una toma de fondo tipo tirolesa podría reducir algunas obras en un cauce rocoso, pero necesitaría evaluar sedimentos, obstrucción, pendiente y estabilidad del lecho. No se selecciona una geometría definitiva de azud sin sección del río ni estudio de avenidas.

Para velocidad de aproximación 0.6 m/s, porosidad 0.60 y fracción libre de obstrucción 0.60, el área bruta ilustrativa es:
\begin{equation}A_{reja}\geq\frac{1.90}{0.60\times0.60\times0.60}=\dato{RejaArea}\ \mathrm{m^2}.\end{equation}
La distribución de esa área entre paneles y su profundidad se resolverá con niveles reales del río. No se afirma que exista una lámina suficiente para una reja de dimensiones arbitrarias. El espaciamiento de barras, su resistencia, el mecanismo de limpieza y la velocidad local se especificarán con el proveedor de turbina y el estudio sedimentológico.

\section{Canal y desarenador}
Se propone canal rectangular revestido, ancho 1.60 m, altura interior 1.30 m, longitud 150 m, pendiente 0.002 y Manning \(n=0.015\). Se resuelve iterativamente:
\begin{equation}Q=\frac{1}{n}A R_h^{2/3}S^{1/2},\quad A=by,\quad R_h=\frac{by}{b+2y}.\end{equation}
A 1.90 \(\mathrm{m^3/s}\), el tirante es \dato{CanalY} m, la velocidad \dato{CanalV} m/s y el número de Froude \dato{CanalFr}. El borde libre resultante es \dato{CanalFree} m. La capacidad a tirante 1 m es \dato{CanalCap} \(\mathrm{m^3/s}\). Deben verificarse transiciones, curvas, resalto, juntas, filtraciones y velocidades erosivas con el perfil real. La altura no se utiliza para transportar una avenida del río.

Se estudian dos cámaras de desarenación de 29 por 4 m y tirante 1.80 m. Cada cámara puede tratar el caudal completo cuando la otra se limpia. Con velocidad de caída supuesta \(w_s=0.025\) m/s y factor 1.5:
\begin{equation}A_s\geq1.5\frac{Q_d}{w_s}=114\ \mathrm{m^2};\qquad A_{adoptada}=116\ \mathrm{m^2}\text{ por cámara}.\end{equation}
Este criterio de sedimentación ideal debe comprobarse con granulometría, temperatura, concentración, turbulencia y longitud de entrada. No garantiza eliminación de toda partícula abrasiva. La purga se devuelve de forma controlada sin provocar erosión ni afectar el tramo ecológico.

\section{Cámara de carga, alivio y restitución}
La cámara tiene planta 8 por 5 m y profundidad útil de operación 2.5 m, con volumen geométrico de 100 \(\mathrm{m^3}\). Se consideran 60 \(\mathrm{m^3}\) aprovechables entre niveles de regulación, equivalentes a \dato{CamSeconds} s al caudal máximo. No es un embalse de regulación diaria. La cota de operación, sumergencia de entrada, vórtices y volumen de maniobra necesitan una geometría definitiva.

El alivio se verifica para el caudal derivado completo, con \(Q=C b h^{3/2}\), \(C=1.7\) y \(h=0.5\) m; requiere un ancho aproximado \dato{VertAncho} m. Este alivio protege la conducción ante rechazo de carga; \textbf{no sustituye el paso de avenida en el río}. Debe conducir a una descarga resistente a la erosión, independiente del acceso y de las fundaciones.

La restitución debe ubicarse aguas abajo de la casa de máquinas, manteniendo descarga libre o la contrapresión admitida por la turbina. Se requiere perfil del río, niveles de avenida y verificación de socavación. No se asigna una coordenada definitiva donde no hay identificación de campo.
\begin{figure}[htbp]\centering\includegraphics[width=.96\linewidth]{figuras/obras-r03.png}\caption{Secuencia de obras y dimensiones de estudio. Esquema sin escala; no contiene secciones constructivas certificadas.}\end{figure}

\section{Pérdidas y punto nominal}
Para acero de rugosidad supuesta 0.045 mm se calcula Darcy--Weisbach y el factor de fricción con una aproximación explícita de Colebrook, usando \(\nu=1.24\times10^{-6}\ \mathrm{m^2/s}\). Las pérdidas localizadas suman \(K=1.2\): entrada 0.5, dos cambios de dirección 0.4 y válvula/transición 0.3. No incluyen un inyector como si fuese una pérdida ajena a la eficiencia de turbina.
\begin{equation}h_f=f\frac{L}{D}\frac{v^2}{2g},\quad h_l=\sum K\frac{v^2}{2g},\quad H_n=H_b-h_f-h_l.\end{equation}
\input{generado/diametros}
La alternativa principal usa diámetro interior nominal \dato{Di} mm. Para producir 600 kW en bornes, el caudal calculado es \dato{Qnom} \(\mathrm{m^3/s}\), inferior al límite de 1.90. Resultan \(v=\dato{Vel}\) m/s, \(Re=\dato{Re}\), \(f=\dato{Fdarcy}\), \(h_f=\dato{Hf}\) m y \(h_l=\dato{Hloc}\) m, con salto neto \dato{Hn} m. Si se entregara todo el caudal máximo, la potencia hidráulica excedente no autoriza sobrecargar los generadores: se regula la admisión y se descarga el sobrante.

\chapter{Turbina, generador y regulación aislada}
\section{Comparación de tecnologías}
La Michell--Banki es un antecedente referido por el usuario y por la memoria de línea, sin placa recuperada. Se compara como alternativa de diseño. La documentación del fabricante OSSBERGER indica un campo amplio de salto y caudal que incluye el caso estudiado, así como admisión dividida (H12). Ello confirma aplicabilidad general, no eficiencia garantizada de un equipo a medida.
\begin{longtable}{L{2.3cm}L{4.2cm}L{7cm}}\toprule Tecnología & Adecuación al caso & Aspectos de selección \\ \midrule
Michell--Banki & Salto medio y caudal variable & Construcción y mantenimiento accesibles; admisión dividida y modularidad. Rendimiento pico usualmente menor que una Francis optimizada; exige control de abrasión.\\
Francis & Puede operar en el rango & Mayor rendimiento pico potencial; selección depende del número específico de revoluciones, cavitación, altura de instalación y carga parcial. Más exigente en mantenimiento y regulación.\\
Pelton / Turgo & Evaluables técnicamente & Turbinas de impulso; a este salto y caudal pueden requerir varios chorros y equipos mayores. El antecedente Pelton de Phara tiene un salto muy distinto y no decide esta selección.\\ \bottomrule\end{longtable}
Se propone estudiar tres Banki de 200 kW eléctricos, sin afirmar que la central existente tenga tres unidades. La modularidad permite apagar grupos en estiaje y mantener otros durante mantenimiento. La alternativa de dos grupos de 200 kW conserva interés como etapa si se confirma menor demanda. La elección final debe comparar cotizaciones y curvas con igual salto, caudal, garantía, repuestos y costo de ciclo de vida.

\section{Punto nominal y generador}
Cada grupo se propone con generador síncrono de 250 kVA, 200 kW a factor de potencia nominal 0.8, 380 V, 60 Hz y ocho polos:
\begin{equation}n_s=\frac{120f}{p}=900\ \mathrm{rpm}.\end{equation}
La velocidad de la turbina debe ser compatible con el generador o incluir transmisión dimensionada. Una estimación preliminar de velocidad periférica \(u\simeq0.48\sqrt{2gH_n}\) conduce a diámetro de rodete cercano a \dato{Rodete} mm; un chorro de ancho supuesto 0.75 m tendría espesor aproximado \dato{JetWidth} mm. Son relaciones de predimensionamiento, no un diseño de álabes. El torque por grupo al punto nominal es aproximadamente \dato{Torque} N m.

Se adoptan eficiencias nominales \(\eta_t=0.84\), \(\eta_m=0.985\) y \(\eta_g=0.925\). El balance da \dato{Ph} kW hidráulicos, \dato{Pt} kW a la salida de turbinas, \dato{Pacople} kW de acople y \dato{PG} kW eléctricos. Las curvas parciales y el caudal mínimo por unidad son hipótesis del modelo y deben sustituirse por datos garantizados. Los 250 kVA deben ser \textbf{aptos para la altitud y temperatura del sitio}; no se aplica un porcentaje de desclasificación inventado. IEC 60034-1:2026 se consulta como referencia de especificación, con acceso al catálogo y no al texto completo (N11).

\section{Frecuencia, tensión y reparto entre grupos}
Una red aislada requiere que generación y carga se equilibren instantáneamente. Se propone un grupo líder de frecuencia y los restantes con estatismo coordinado, o un controlador de reparto equivalente. El regulador hidráulico ajusta el caudal a medio plazo; un controlador electrónico de carga (ELC) absorbe rechazos transitorios con cargas de disipación. El AVR mantiene tensión con reparto de reactivos; no deben competir tres controles isócronos independientes sobre la misma barra.

Se prevén tres módulos de carga de disipación de 200 kW, con capacidad transitoria total de 600 kW. Un rechazo completo durante ocho segundos disiparía aproximadamente 1.33 kWh. El sistema de refrigeración, protección térmica y descarga debe soportar esa potencia instantánea; la pequeña energía no permite minimizar el equipo. No se atribuye generación útil a la energía disipada. El gobierno reduce admisión con cierre controlado de al menos ocho segundos, sujeto a validar el transitorio con geometría y proveedor.

Con constante de inercia supuesta de 2 s, un rechazo total puede producir una variación inicial de frecuencia del orden de 15 Hz/s. Una respuesta electrónica del orden de un ciclo es un objetivo de diseño, no una prestación comprobada. Deben ensayarse arranque, incorporación y retiro de grupos, saltos de carga, falla del ELC, falla de AVR y pérdida de respaldo. La protección de sobrevelocidad debe disponer de actuación independiente del regulador, conforme a CNE Suministro 130.
Las funciones de regulación y disparo se representan en la lámina PL-07, junto con los criterios de prueba y protección.

\Needspace{24\baselineskip}
\section{Servicios auxiliares y arranque autónomo}
\begingroup\small\input{generado/auxiliares}\endgroup
A plena carga se requieren \dato{Paux} kW y \dato{Qaux} kvar de auxiliares. La ventilación considera las pérdidas de turbinas y generadores dentro de la casa, auxiliares y cableado: aproximadamente \dato{Heat} kW térmicos. Con densidad de aire supuesta 0.9 \(\mathrm{kg/m^3}\), \(c_p=1005\ \mathrm{J/(kg\,K)}\) y aumento de 10 K, se necesitan \dato{AirFlow} \(\mathrm{m^3/s}\). A 100 Pa y rendimiento 0.6, la potencia de ventiladores es del orden de \dato{FanPower} kW. Se proponen dos ventiladores de 1.1 kW; el caudal real debe comprobarse con su curva y pérdidas de rejillas. El transformador principal se ubica en exterior.

El transformador auxiliar es de 20 kVA, 380/380--220 V, con neutro secundario propio. No se confunde 220 V entre fases con 220 V fase--neutro. Se propone batería de 24 V y 100 Ah para control, disparo y señalización; a 0.5 kW, profundidad de descarga 0.8 y eficiencia 0.9 proporciona \dato{DCtime} h de autonomía teórica. No alimenta motores de gobernadores en corriente alterna. El arranque autónomo necesita un módulo diésel, mando hidráulico manual o acumulador dimensionado; debe definirse la secuencia de energización de auxiliares antes de aceptar carga.

\section{Respaldo y límites operativos}
Para cubrir el escenario conservador durante períodos sin caudal se estudian tres grupos diésel de 200 kW. Su capacidad en usuarios es aproximadamente \dato{PBackupUsers} kW bajo el mismo sistema de pérdidas. La potencia combinada hidro--diésel se limita a 600 kW en la barra actual, mediante control e interbloqueo: no se diseña una barra de 1200 kW de exportación. La operación diésel requiere combustible almacenado con contención, ventilación, protección contra incendios, escape, control de ruido y mantenimiento; ni la disponibilidad ni el suministro de combustible están garantizados. El costo de este respaldo cuestiona la viabilidad y se analiza en el capítulo económico.

\chapter{Tubería de presión, transitorios y estructuras}
\section{Material, diámetro y presión estática}
Se propone tubería de acero soldado, diámetro interior nominal de 1000 mm y espesor nominal de 16 mm, con recubrimiento exterior e interior y protección según agresividad del medio. El diámetro exterior nominal es 1032 mm. Se especifica un límite de fluencia mínimo de 250 MPa como requisito de compra, sin atribuirlo a una tubería existente. Se consideran tolerancia negativa de espesor de 12.5\% y pérdida futura por corrosión de 1 mm: espesor resistente al final de vida de \dato{EspEf} mm. Estos márgenes son hipótesis que deberán sustituirse por especificación de fabricación y evaluación de corrosión.

La presión estática asociada al salto de 45.3 m es aproximadamente 4.44 bar manométricos en el extremo bajo. En régimen, las pérdidas reducen la presión disponible, pero el cierre puede generar una sobrepresión mucho mayor. Una clase nominal no se selecciona únicamente a partir del salto estático. Alternativas en material plástico requerirían verificar temperatura, fluencia, uniones, anclajes, instalación, golpe de ariete y disponibilidad comercial a gran diámetro; no se les atribuye resistencia equivalente al acero sin cálculo.

\section{Golpe de ariete y límites del modelo}
Se desarrolla un modelo unidimensional por el método de las características (MOC): embalse de carga de nivel fijo, tubería uniforme, fricción Darcy, y válvula equivalente con \(Q=k\tau\sqrt{H}\). Se usa paso espacial uniforme con 100 segmentos y paso temporal que cumple Courant 1. La celeridad se calcula con compresibilidad del agua y elasticidad del acero. Se analizan cierres de 0, 1, 4 y 8 s tanto para pared nueva mínima de 14 mm como para pared final de 13 mm. El diámetro interior aumenta coherentemente al reducir pared; se incrementa también rugosidad a 0.15 mm al final de vida.

El caudal inicial es 1.90 \(\mathrm{m^3/s}\), mayor que el caudal nominal de operación, para estudiar el máximo hidráulico. Se incluye un perfil ideal de descenso continuo; los puntos altos reales y el cierre efectivo de una Banki pueden modificar las ondas. La relación de Joukowsky \(\Delta H=a\Delta v/g\) se utiliza como comprobación del primer salto para un cierre instantáneo sin fricción, no como sustituto del análisis temporal.
\input{generado/moc}
Se detiene la integración al detectar presión manométrica inferior a una referencia de --6.8 m de columna de agua, representativa del límite atmosférico--vapor a la altitud del sitio. El valor es un supuesto conservador, no una medición barométrica. Después de ese instante el modelo monofásico pierde validez: no se publican presiones algebraicas físicamente imposibles ni se certifica el pico posterior al colapso de cavidades. El cierre rápido presenta este problema.

Para cierre controlado de ocho segundos se obtiene una máxima de estudio cercana a \dato{PNslow} bar. La primera máxima válida de cierre brusco alcanza aproximadamente \dato{PNfast} bar. Se estudian válvulas y accesorios \textbf{PN40 como candidatos}, pero esa clase no demuestra por sí sola seguridad ante separación de columna, vacío, fatiga o colapso de cavidades. Se exige validar con el perfil real, curva de cierre, condiciones de admisión, regulación y modelo bifásico cuando corresponda. El ELC debe tomar el rechazo rápido mientras el gobierno cierra lentamente. Debe evitarse que una orden eléctrica provoque cierre hidráulico instantáneo.
\begin{figure}[htbp]\centering\includegraphics[width=.95\linewidth]{figuras/transitorio-r03.png}\caption{Resultados válidos del MOC y envolvente del cierre controlado. Las trazas truncadas indican pérdida de validez por presión baja.}\end{figure}

\section{Espesor resistente y comprobación de membrana}
Se evalúa esfuerzo circunferencial \(\sigma_\theta=pD/(2t_e)\), longitudinal de membrana \(\sigma_z\simeq\sigma_\theta/2\) y equivalente de von Mises. Con el estado de pared y la presión compatibles, resultan \(\sigma_\theta=\dato{SigmaHoop}\) MPa y \(\sigma_{VM}=\dato{SigmaVM}\) MPa. El admisible de estudio es \(250\times0.85/1.5=\dato{SigmaAllow}\) MPa, usando eficiencia de junta 0.85 y factor 1.5 \textbf{supuestos}, no valores atribuidos a una cláusula no consultada.

La comprobación de membrana es favorable bajo esos supuestos, pero no cubre flexión local de silletas, ovalización, pandeo por vacío, concentraciones en derivaciones, fatiga ni estabilidad sísmica. ASCE MOP 79, segunda edición 2012 (N26), identifica criterios complementarios para tuberías de acero; se consultó su ficha y estructura, no el texto completo. USBR EM03 (1966, N20) se utiliza como referencia histórica de mecanismos, sin adoptar códigos antiguos como obligación peruana.

\section{Apoyos, anclajes y acciones térmicas}
Se comparan vanos de 4, 6 y 8 m mediante una viga ideal bajo peso de acero y agua. Para metrados se usa separación preliminar de 4 m, que produce 89 apoyos incluyendo extremos. La masa nominal de acero es aproximadamente \dato{MassPipe} t antes de accesorios. El cálculo global no dimensiona una silleta: deben verificarse contacto, ángulo de apoyo, refuerzos anulares, reacciones, fundación, deslizamiento y posibilidad de movimiento térmico.
\input{generado/apoyos}

Para \(\Delta T=40\ ^\circ\mathrm{C}\), la expansión libre de los 350 m es aproximadamente \dato{Expansion} mm. Se requiere una disposición coordinada de puntos fijos, apoyos deslizantes y juntas, sin permitir fuerzas térmicas incompatibles con los anclajes. En válvulas, codos y cambios de pendiente se evalúan empujes y fuerzas dinámicas; los anclajes no pueden ser simples dados sin verificación.

Una equivalencia de peso--fricción para un codo de estudio de 45 grados, usando \(\mu=0.6\), peso específico de concreto de 24 \(\mathrm{kN/m^3}\) y factor 1.5, arroja un volumen de gravedad cercano a \dato{AnchorVol} \(\mathrm{m^3}\) ante el primer pico brusco válido. Este resultado muestra la magnitud del problema; \textbf{no es un bloque de concreto diseñado}. Se considera estudiar anclaje a roca competente, con caracterización geotécnica, extracción, arrancamiento, estabilidad global y pruebas. Se presupuestan conjuntos provisionales y no un volumen ejecutivo inventado.

\section{Estructuras, geotecnia y sismo}
El diseño académico adopta la E.030 modificada por RM 183-2026-VIVIENDA y considera la disposición transitoria de RM 217-2026. El anexo distrital sitúa Limbani en zona 2, \(Z=0.25\). La categoría A2 comprende instalaciones de generación y transformación, con \(U=1.5\). No se pueden completar espectro, fuerzas y derivas sin perfil de suelo, sistema estructural, períodos y estudio de mecánica de suelos. E.050 (2018) guía la necesidad y contenido del estudio geotécnico; E.060 y demás normas estructurales aplicables se verifican antes de diseñar armado y fundaciones. El acero de refuerzo usado en presupuesto es un indicador de metrado, no un cálculo estructural.

Se proponen válvula principal de aislamiento, órganos de admisión por grupo, purgas, drenes, medición de presión y dispositivos de aire donde el perfil lo requiera. Su tamaño, tipo, capacidad y tiempos de actuación deben coordinarse con el MOC. No se adopta una válvula de aire genérica como solución automática al golpe de ariete. El plan de inspección incluirá soldaduras, espesores, recubrimiento, estanqueidad, instrumentación y pruebas de presión definidas por la especificación aplicable.

\chapter{Diseño eléctrico y seguridad del sistema aislado}
\section{Arquitectura, transformación y límites}
La generación opera inicialmente aislada del SEIN. Tres generadores de 250 kVA se conectan a 380 V, con interruptor por grupo y barra principal; el respaldo diésel tiene tres módulos de 200 kW. Se limita la potencia total simultánea de generación a 600 kW. Un transformador de 800 kVA, 10/0.38 kV, eleva hacia una red primaria académica de 8 km y seis subestaciones de distribución: tres de 160 kVA, dos de 100 y una de 50 kVA, total 730 kVA. Las potencias de placa de distribución no se suman como demanda coincidente.

El grupo de conexión propuesto mantiene referencia AT/BT: devanado de 10 kV en delta y devanado de 380 V en estrella con neutro accesible, con desplazamiento horario por especificar con el fabricante. Invertir el flujo de potencia no permite renombrar arbitrariamente el grupo vectorial. El transformador auxiliar de 20 kVA es 380/380--220 V, con estrella secundaria y neutro propio para utilización.
\begin{figure}[htbp]\centering\includegraphics[width=.98\linewidth]{figuras/unifilar-r03.png}\caption{Unifilar académico del sistema aislado. Sin fuente SEIN; impedancias, interruptores y neutros requieren coordinación definitiva.}\end{figure}
\input{generado/balance}
Las pérdidas del cuadro se descuentan una sola vez, usando factor de potencia de transporte 0.90 como hipótesis conservadora; la demanda doméstica conserva 1.0 de N21. Resultan \dato{Pusers} kW en usuarios y margen de \dato{Margin} kW sobre el pico. Arranques, temperatura, desequilibrio o caída adicional pueden consumir ese margen. El balance no demuestra cumplimiento de calidad.

\section{Conductores y caída de tensión}
La corriente nominal por generador es \dato{Iunit} A y la suma de placas, \dato{Ibus} A. El transformador de 800 kVA tiene corriente nominal BT de \dato{Itrafo} A. Se estudia interruptor de 630 A por grupo y de 1600 A en barra, con capacidad de interrupción candidata de 25 kA. Para equipos marcados al 80\% de carga continua, 630 y 1600 A permiten 504 y 1280 A respectivamente. La selección final debe considerar temperatura, fabricante, modo de instalación y coordinación.

La regla CNE Utilización 050-104 distingue corriente continua, marcado del equipo y tablas de ampacidad. Las condiciones de las tablas de enterramiento D8A/D8B son específicas: no corresponde multiplicar indiscriminadamente toda ampacidad por 0.8. Como referencia, para cobre de 240 \(\mathrm{mm^2}\) el detalle 4 consultado indica 446 A en D8A y 442 A en D8B, y otros detalles dan valores diferentes. Se proponen dos conductores de cobre de 240 \(\mathrm{mm^2}\) por fase y generador, longitud 25 m. Cada uno conduce aproximadamente 190 A, pero su ampacidad admisible debe verificarse con disposición, agrupamiento, temperatura, resistividad térmica y reparto entre paralelos. Se requiere igualdad de material, sección y longitud, conexiones equivalentes y PE dimensionado según el sistema.

La resistencia de estudio por conductor es 0.095 ohm/km; la caída del ramal de generación es aproximadamente \dato{DropCable}\%. Los 450 m presupuestados corresponden a \(3\) grupos por \(3\) fases por \(2\) paralelos por \(25\) m; no incluyen ficticiamente un neutro de potencia en cada fase. Las conexiones de respaldo se consideran con cuatro paralelos de cobre de 400 \(\mathrm{mm^2}\) por fase a lo largo de 25 m, sujeto a revisión térmica.

Para la línea primaria se propone AAAC de 95 \(\mathrm{mm^2}\), resistencia a 20 grados de 0.35 ohm/km y coeficiente térmico 0.0036/\(\!^\circ\mathrm{C}\). A 40 grados resulta aproximadamente 0.3752 ohm/km. La reactancia 0.38 ohm/km es una hipótesis de disposición; deben fijarse geometría y flechas. A plena carga la caída es \dato{DropMT}\%. El recorrido de 8 km es un corredor aproximado para cubrir las localidades, no una línea levantada ni un trazado aprobado. El diseño mecánico de postes, vanos y distancias CNE debe desarrollarse con topografía real.

Se evalúa una rama BT de 25 kW, 150 m, 380 V y conductor autoportante de aluminio 3 por 35 más 16/25 \(\mathrm{mm^2}\), según cuadro de N21: \(R_{40}=0.930\) ohm/km, \(X=0.103\) ohm/km y corriente de referencia 129 A. Resulta corriente aproximada de 38 A y caída \dato{DropBT}\%. La norma fija sección mínima de fases de 25 \(\mathrm{mm^2}\); no se arrastra la antigua selección de 16 como conductor general de fases. Su numeral 4.5 contempla 7.5\% máximo BT para SER bajo su ámbito; ello no reemplaza los criterios de utilización ni demuestra tensión correcta en todos los usuarios. La regulación del transformador, reparto por fase, longitud real, arranque y calidad deben analizarse conjuntamente. La pérdida total BT de 3\% es una hipótesis del balance, no el resultado de todos los circuitos individualizados.

\section{Cortocircuito y protección de una red débil}
Se utiliza reactancia subtransitoria supuesta de 0.20 p.u. para cada generador hidráulico, 0.15 p.u. para respaldo y \(u_k=6\%\) para T1. Se calcula \(I''\simeq I_n/X''\) y se combina la impedancia equivalente con el transformador. El resultado inicial simétrico es \dato{IKH} kA para hidro, \dato{IKB} kA para respaldo y una cota simplificada de \dato{IKmax} kA para fuentes simultáneas. La contribución hidroeléctrica referida al lado 10 kV es \dato{IKMT} A. Estos datos no son valores de placa; se requiere impedancia de secuencia, excitación, motores, fuente diésel, posición y tipo de falla.

IEC 60909-0:2026, edición 3 (N25), sustituyó la edición de 2016 según la ficha oficial. No se obtuvo el texto completo, por lo que no se declara que este cálculo simplificado cumpla íntegramente el método IEC. El Icu candidato de 25 kA debe revisarse con corriente de cierre, componente continua, relación X/R, capacidad térmica y dinámicas. Los cortocircuitos monofásicos y extremos de red no pueden inferirse del máximo trifásico de barra.

La corriente de falla de un generador disminuye con el tiempo: con \(X_d=1.8\) p.u. supuesto, el agregado hidráulico puede aproximarse a 633 A sostenidos, inferior a la corriente de carga plena. \textbf{Una protección de sobrecorriente 50/51 por sí sola puede no detectar una falla remota}. Se prevén protección diferencial de generador y transformador (87G/87T), sobrecorriente restringida o controlada por tensión (51V), fallas a tierra de acuerdo con neutros (64G/59N/67N), potencia inversa (32), pérdida de excitación (40), secuencia negativa (46), tensión y frecuencia (27/59/81), y comprobación de sincronismo (25) entre grupos internos. Los ajustes necesitan curvas de decremento, capacidades térmicas, transformadores de medida y estudio de selectividad; no se publican ajustes inventados.
\input{generado/fallas_remotas}
El cuadro agrega impedancia de 8 km de línea, transformación SED con \(u_k=4\%\) supuesto y rama BT de 150 m. Se compara operación de tres grupos, un solo grupo y un límite sostenido de un grupo, sin asumir que este último sea una falla estacionaria garantizada con excitación real. La caída de corriente confirma la necesidad de sensibilidad y selectividad. No se derivan de este cuadro corrientes fase--tierra, impedancia homopolar ni ajustes definitivos.

\section{Neutros y puesta a tierra}
Se propone estudiar puesta a tierra del neutro de generación mediante impedancia coordinada y un único esquema activo compatible con grupos en paralelo. La delta de 10 kV no aporta un neutro: si se adopta puesta a tierra de esa red por impedancia será necesario transformador zigzag y resistencia dimensionados. Las subestaciones de distribución y el auxiliar tendrán sus neutros BT y conductores de protección definidos según el sistema de utilización. No se conectan resistencias de neutro como si fueran conductores PEN de viviendas.

La malla conceptual de casa de máquinas es 20 por 12 m, paso 4 m, cobre de 50 \(\mathrm{mm^2}\), longitud horizontal 152 m y ocho electrodos de 2.4 m. La sección no se certifica sin corriente a tierra y duración de despeje. El estimador de gabinete \(R_g\simeq\rho/(4\sqrt{A/\pi})+\rho/L\) sirve para sensibilidad, sin acreditar el aporte de varillas ni efectos de capas del suelo.
\input{generado/tierra}
El GPR es elevación de potencial de la malla y \textbf{no equivale a tensión de contacto}. CNE Suministro 036 exige considerar paso, contacto y transferencia; su referencia de 25 ohm para un electrodo individual no es un criterio global suficiente de seguridad. Tampoco existe aquí una resistencia universal de 5 ohm que cierre la verificación. Se requiere medir resistividad por estratos, definir reparto de corriente, tiempo de despeje, superficie y cercos, y calcular tensiones accesibles con metodología reconocida. La ficha IEEE 80-2013 consultada indica condición reservada/inactiva; se debe seleccionar y adquirir una edición contractual reconocida, sin presentarla como texto completo vigente consultado.
\begin{figure}[htbp]\centering\includegraphics[width=.95\linewidth]{figuras/puesta-tierra-r03.png}\caption{Malla y separación conceptual de funciones de neutro. Dimensiones de estudio; verificación de paso/contacto pendiente.}\end{figure}

\section{Espacios y seguridad}
CNE Suministro 125 y tablas correspondientes definen espacios según tensión y condición de exposición. Para BT 151--600 V, los casos de la tabla 125-1 consultada exigen 900, 1070 o 1200 mm; no se aplica una sola distancia a cualquier tablero. La planta académica reserva 1.5 m frontal, y la media tensión debe contrastarse con su tabla aplicable y equipo real. Se requieren segregación, acceso restringido, barreras, señalización, bloqueo y consignación, extinción compatible, control de incendio de combustible, salidas y plan de rescate. Distancias aéreas de suministro y zonas de trabajo se verifican en el proyecto de distribución, considerando relieve y cruces reales.

\section{Adaptación eventual al SEIN entre 2046 y 2076}
La memoria oficial de Electro Puno 2023 incluye transformación 10/22.9 kV y redes regionales de 10, 22.9 y 60 kV (H14). No acredita un punto disponible ni una conexión aprobada en Limbani. Se adopta 10 kV en el diseño actual por el antecedente local; no se elige automáticamente 22.9 kV para simular una futura interconexión.

Como previsión se reserva espacio para una celda adicional, canalizaciones de comunicaciones, medición y acceso a transformación exterior. Puede estudiarse equipamiento con aislamiento de clase 24 kV cuando su costo y funcionamiento a 10 kV resulten adecuados, pero elevar la tensión futura exigiría verificar o sustituir transformadores, aislamiento, conductores, distancias y protecciones. Las decisiones se tomarán con la red y normativa existentes en aquel momento.

Una interconexión requeriría estudio de conexión y autorización de la concesionaria, flujos de carga, cortocircuito, estabilidad, sincronismo, puesta a tierra, calidad, protección de interfaz y antiisla, medición comercial, comunicaciones, actualización ambiental, servidumbres y concesiones pertinentes. El horizonte 20--50 años es orientativo; no altera la operación aislada ni se utiliza como garantía de crecimiento o rentabilidad.

\chapter{Normativa y ruta de autorizaciones}
\section{Aplicación de CNE Suministro y Utilización}
El CNE Suministro 2011 regula los aspectos de suministro, estaciones, líneas, conductores, distancias, maniobra y seguridad de la infraestructura. El CNE Utilización, aprobado por RM 037-2006 con la modificación RM 175-2008, se aplica a las instalaciones de utilización que corresponden, incluidos los circuitos auxiliares y receptores. No se emplea una regla residencial para sustituir el diseño de una estación generadora ni se atribuye al CNE una potencia por vivienda que procede de la norma rural de 2023.

La matriz del anexo A registra disposición, criterio, evidencia y estado abierto. Se revisaron particularmente CNE Suministro 036 (puesta a tierra), 125 (espacios), 130 (sobrevelocidad), 132 (neutros) y reglas de líneas 230, 232 y 234. CNE Utilización 030-004, 050-104, 050-106, sección 060 y tablas de ampacidad se consultan junto a instalación real y marcado de equipos. La existencia de una referencia bibliográfica no acredita cumplimiento de todas sus reglas.

\section{Cambios normativos importantes para esta revisión}
La RD 108-2023-MINEM/DGE sustituye las antiguas bases rurales de redes secundarias aprobadas por RD 031-2003. Sus valores de calificación y simultaneidad modifican sustancialmente la demanda, por lo que se recalculan ambos y se conserva el ámbito académico de aplicación. La E.030 fue modificada en 2026; no se reutiliza una zonificación histórica sin verificar el anexo distrital actual. Se identifica además la disposición transitoria RM 217-2026, cuya aplicación a una obra real depende de las condiciones de tramitación.

La ficha oficial de IEC 60909-0 identifica una edición 2026 que reemplaza la de 2016. IEC 60034-1:2026 y IEC 60076-1:2011 sirven como referencias de especificación de máquinas y transformadores. La consulta disponible de estas normas internacionales es de catálogo, no de sus textos completos pagados. No se inventan subcláusulas ni se certifican diseños conforme a ellas. Antes de una contratación se deben adquirir las ediciones pertinentes y precisar ensayos, pérdidas, elevación de temperatura, cortocircuito, aislamiento y funcionamiento a la altitud del sitio.

\section{Agua, ambiente, estructuras y seguridad}
Ley 29338 y su Reglamento distinguen disponibilidad, derecho de uso, intervención en cauce y autorización de obras. Se revisan, entre otros, artículos 44, 47 y 53 de la ley para uso y licencia; 74 para faja marginal; 81 para evaluación ambiental; y 104 para obras. La modificación DS 023-2014-MINAGRI se consulta junto a los artículos reglamentarios pertinentes. El procedimiento ANA de una ejecución real deberá verificarse con el TUPA y las disposiciones vigentes; no se declara lectura integral de anexos no recuperados de RJ 007-2015.

El caudal ecológico se desarrolla según RJ 267-2019, numeral 6: metodología dependiente de categoría y coordinación con la autoridad. El porcentaje de estudio utilizado no es una obligación legal ni un caudal aprobado. Deben examinarse usos poblacionales, agrarios, energéticos y restricciones del tramo, así como disponibilidad en estiaje y retorno del agua.

El DS 014-2019-EM, junto a su fe de erratas, regula protección ambiental de actividades eléctricas. La categoría no se determina solamente porque la potencia sea pequeña. El tipo de aprovechamiento, embalse, longitud de línea, ubicación y sensibilidad ambiental pueden modificar el instrumento exigible. Se revisan criterios del anexo correspondiente sin afirmar una DIA automática. El DS 011-2026-EM se refiere a adecuación de actividades en curso; no constituye certificación ambiental de una nueva ampliación. Debe distinguirse situación legal de la instalación existente, rehabilitación y obras nuevas.

Para obras civiles se aplicará el RNE pertinente: E.030 para sismo, E.050 para suelos y cimentaciones, E.060 para concreto armado, además de reglas de arquitectura, accesos y seguridad que correspondan. El diseño estructural definitivo requiere geotecnia y modelos de carga. Ley 29783 y el reglamento de seguridad eléctrica RM 111-2013 sustentan gestión de riesgos, competencias, consignación, trabajo en altura y planes de emergencia. No se presenta esta memoria estudiantil como documentación firmada por profesionales habilitados para ejecutar.

\section{Obligaciones y recomendaciones técnicas}
Las leyes, reglamentos y códigos nacionales aplicables son obligaciones dentro de su ámbito. ASCE MOP 79, USBR EM03, fichas de fabricante y métodos hidráulicos son apoyos técnicos; solo una especificación contractual o incorporación normativa puede atribuirles otra exigibilidad. Se separan parámetros adoptados de valores medidos y se identifican las ediciones. Donde falta acceso al texto de una norma complementaria, se deja abierta la comprobación correspondiente.

\section{Ruta orientativa para una ejecución real}
Todos los procedimientos están pendientes; no hay permisos obtenidos.
\begin{enumerate}
\item Identificar titular, estado legal y bienes de la central existente. Levantar instalaciones y verificar propiedad, acceso, servidumbres y ámbito de prestación. La mención Municipalidad/Junta en antecedentes no acredita titularidad.
\item Consultar con ANA/AAA competente disponibilidad, derechos concurrentes, caudal ecológico, faja marginal y autorizaciones necesarias. Tramitar los procedimientos aplicables y la licencia de uso; no atribuir la licencia de Phara a Limbani.
\item Definir con la autoridad ambiental el instrumento y la categoría para rehabilitación, ampliación y distribución. Incorporar línea, combustible, sedimentos, ruido, afectación de cauce y patrimonio. Evaluar permisos forestales, arqueológicos y prediales según afectaciones reales.
\item Verificar competencia MINEM/Gobierno Regional de Puno y el régimen eléctrico vigente. La modificación introducida por DL 1002 considera concesión definitiva para generación hidráulica por encima de 500 kW. El caso estudiado de 600 kW supera ese umbral. También debe evaluarse la distribución de servicio público con demanda superior a 500 kW y la autorización de generación térmica de respaldo superior a 500 kW conforme a los artículos aplicables de LCE. No se conserva una pretendida exención general basada en los 150 kW del borrador.
\item Desarrollar estudios de ingeniería firmados, expediente de obras, autorizaciones de intervención, permisos municipales y de seguridad que procedan. Verificar la condición de sistema eléctrico rural y la normativa específica con la autoridad responsable.
\item Obtener las autorizaciones, ejecutar con control de calidad y realizar pruebas de presión, regulación, protecciones, aislamiento y tierra antes de puesta en servicio. Registrar operación, mantenimiento y responsables.
\end{enumerate}

\chapter{Metrados, presupuesto y viabilidad}
\section{Herramientas y fronteras del metrado}
El presupuesto se genera desde parámetros canónicos mediante Python, con salida CSV editable. Se evita transcribir totales de capítulos distintos. Las longitudes, equipos, usuarios, auxiliares y potencias son las mismas del modelo de diseño. Se conserva la memoria de cálculo y el manifiesto de fuentes; cualquier modificación de diámetro, grupos o cobertura obliga a reconstruir cantidades, balance y costos.

Para obras civiles se estiman espesores y sobreanchos académicos. El canal produce 101.25 \(\mathrm{m^3}\) de concreto, las dos cámaras de desarenación 143.43 y la cámara de carga 34.28, total aproximado 278.95. El refuerzo se estima en 100 \(\mathrm{kg/m^3}\); no representa planos de armado. La casa de 240 \(\mathrm{m^2}\) se valora por superficie construida, evitando sumarle de nuevo estructura y concreto que ya incluye su partida. Las cimentaciones de silletas y anclajes son provisionales y requieren geotecnia.

La tubería corresponde a 350 m de acero de diámetro exterior 1032 mm y espesor nominal 16 mm. Se agregan 5\% de masa para accesorios de fabricación como estimación. Hay 89 apoyos, cuatro conjuntos principales de válvulas, tres grupos hidráulicos, tres de respaldo, T1 de 800 kVA y auxiliar de 20 kVA. Las redes incluyen 8 km primarios, 9 km secundarios, seis subestaciones, 1037 conexiones potenciales y 226 puntos de alumbrado. Las partidas de redes instaladas incluyen su alcance de montaje y no se recargan como equipos sueltos.

\section{Precios y fecha de referencia}
Los precios unitarios de las partidas son \textbf{estimaciones académicas en soles a octubre de 2026}, no cotizaciones entregadas en Limbani. Se conserva el código, unidad, cantidad, precio, estado y criterio de cada partida en CSV. La consulta minorista pública recuperada permite contrastar cemento Sol tipo I de 42.5 kg a S/30.90 en la página de Sodimac; el producto Yura HS mostrado por el mismo sitio aparece sin existencias. Esa referencia no prueba disponibilidad, transporte ni precio de concreto puesto en obra.

No se recuperaron ofertas verificables para los grupos, tubería fabricada, transporte o subestaciones del proyecto. Se usan valores redondeados de estudio y se marca su condición estimada. Antes de comparar inversión real se solicitarán cotizaciones con impuestos, moneda, fecha, garantía, transporte, montaje y capacidad a la altitud requerida. No se añade un IGV general a valores cuyas bases de impuestos aún no están desagregadas; el presupuesto debe rehacerse homogéneamente con ofertas netas e impuestos para una contratación.

El informe IRENA 2025 sobre costos de generación en 2024 aporta comparadores internacionales: para pequeña hidro, promedios ponderados 2018--2024 de aproximadamente USD 2024/kW 1873 en China, 2534 en Brasil, 3677 en Europa y 4084 en África. No son cotizaciones peruanas ni incluyen necesariamente la distribución y el respaldo de este escenario. No se convierten mecánicamente a soles de 2026 sin tipo de cambio, inflación y equivalencia de alcance. Su finalidad es identificar qué estudios adicionales hacen falta, no validar una precisión aparente.
\begingroup\small\input{generado/presupuesto}\endgroup
\input{generado/resumen_costos}
El transporte y montaje adicional se aplica al 12\% de suministros de tubería, equipos de generación e instalaciones de central; las redes ya se valoran instaladas. Los estudios y supervisión se estiman al 10\% de la base con logística, y la contingencia al 20\% del subtotal resultante. No sustituyen metrados detallados, programación de obra ni costos de acceso. El total de estudio es aproximadamente \textbf{S/\dato{CapexM} millones}. Se presenta una sensibilidad de \(\pm35\%\), S/\dato{CapexLowM}--\dato{CapexHighM} millones; no es un intervalo probabilístico validado.
\begin{figure}[htbp]\centering\includegraphics[width=.88\linewidth]{figuras/costos-r03.png}\caption{Distribución de costos directos estimados. Los recargos se muestran por separado en el resumen.}\end{figure}

\section{Balance económico transparente}
Se evalúa una cobertura híbrida, con demanda 2026 constante durante 20 años y tasa real supuesta del 10\%. No es una predicción de ventas; evita mezclar una tarifa actual no investigada con ingresos futuros ficticios. La energía hidro servida es \dato{EHidro} MWh/año. La energía disponible pero no absorbida no se cuenta como ingreso. Para cubrir el déficit se estima energía diésel en bornes de \dato{DieselE} MWh/año, con eficiencia de entrega supuesta 0.90.

El consumo específico de combustible se adopta en 0.27 l/kWh y el precio en S/5.50/l, ambos \textbf{hipótesis}; no provienen de una cotización local. Se estiman \dato{DieselLitres} mil litros/año y S/\dato{FuelCostM} millones/año de combustible. El consumo real puede crecer a baja carga y debe verificarse con curvas de cada grupo, fletes y almacenamiento. Se propone usar Facilito de OSINERGMIN y ofertas de abastecimiento para sustituir ese precio. La operación y mantenimiento sin combustible se aproxima al 3\% de inversión, S/\dato{OpexM} millones/año; esta simplificación debe reemplazarse por un plan de activos y reposiciones.
\input{generado/viabilidad}
Los ingresos ensayados no son tarifas aprobadas ni recomendaciones de cobro. El valor de equilibrio simplificado es aproximadamente S/\dato{Lcoe}/kWh, incluyendo inversión anualizada, operación y combustible para cubrir el perfil supuesto. No se presenta como un LCOE bancable: faltan reposiciones, estructura financiera, impuestos consistentes, inflación local y riesgos de disponibilidad.

\section{Interpretación y estrategia por etapas}
La rentabilidad financiera de la solución completa es desfavorable en los ingresos considerados. Aumentar potencia instalada no elimina el déficit de estiaje. El tercer grupo hidráulico cubre mejor el pico pero aprovecha poca energía adicional, por lo que 600 kW no queda demostrado como óptimo económico. El padrón y las mediciones pueden justificar menor capacidad y una estrategia gradual.

Se recomienda estudiar primero rehabilitación y reutilización técnicamente verificable de activos, confirmar usuarios permanentes, gestionar demanda y comparar dos grupos de 200 kW con respaldo y una combinación solar--almacenamiento. No se cuantifica ahorro por reutilización ni generación solar sin inspección y recurso correspondiente. Una evaluación social podría considerar beneficios de continuidad, salud y actividad económica, pero requeriría demanda real, criterios públicos aplicables y financiamiento identificado; no se inventa una subvención que vuelva positivo el resultado.

\chapter{Planos, operación y comprobaciones posteriores}
\section{Índice de planos académicos}
El anexo gráfico independiente está preparado para lectura en A3 horizontal. Los mapas identifican WGS84 / UTM zona 19 sur cuando corresponda, norte, unidades y escala gráfica. Los diagramas funcionales se rotulan sin escala. Las dimensiones de planta son de estudio; no equivalen a coordenadas de replanteo.
\begin{longtable}{L{1.4cm}L{7.8cm}L{4.8cm}}\toprule Lámina & Contenido & Condición \\ \midrule
PL-01 & Ubicación, cuenca y localidades & Cartografía pública y referencias; límites de aforo explícitos.\\
PL-02 & Emplazamiento y perfil hidráulico & CM-A/CM-B diferenciadas; cámara y restitución sin posición confirmada.\\
PL-03 & Secuencia hidráulica y dimensiones & Esquema sin escala; no obra de captación definitiva.\\
PL-04 & Tubería, cierre y soportes & MOC y criterios; anclajes sin detalle ejecutivo.\\
PL-05 & Casa de máquinas & Planta 20 por 12 m; equipos y patio de estudio.\\
PL-06 & Diagrama unifilar & Red aislada 380 V / 10 kV; sin conexión aprobada al SEIN.\\
PL-07 & Regulación y funciones de protección & Bloques sin escala y ajustes pendientes.\\
PL-08 & Malla de puesta a tierra & Planta conceptual; seguridad de contacto/paso pendiente.\\ \bottomrule\end{longtable}

\section{Operación y mantenimiento}
Se propone monitorear caudal derivado y ecológico, niveles de cámara, presión, frecuencia, tensión, potencia activa/reactiva, temperatura y estado de protecciones. El arranque verifica reja y desarenador, posición de válvulas, carga hidráulica, auxiliares, batería, ausencia de personal en zonas de riesgo y habilitación de protecciones. Se acelera el primer grupo, se estabiliza tensión/frecuencia y se incorporan cargas; los grupos adicionales se sincronizan con la barra interna. El respaldo no se conecta sin comprobación de sincronismo o transferencia con enclavamiento adecuado.

Ante rechazo de carga actúa el ELC y se reduce admisión con cierre controlado; una falla hidráulica o sobrevelocidad activa una secuencia independiente y coordinada. El circuito prioritario y los relevadores de frecuencia gestionan déficit, sin dejar salud y comunicaciones a una decisión manual tardía. Las consignas deben ensayarse antes de servicio. Las maniobras con tensión requieren procedimientos aprobados y personal competente; la memoria estudiantil no autoriza esas tareas.

El mantenimiento incluye limpieza segura de rejas, purga alternada del desarenador, inspección de taludes y cauce, control de sedimentos, espesores y corrosión de tubería, juntas, silletas y anclajes. En equipos se revisan rodamientos, vibración, alineación, lubricación, AVR, reguladores, balastos, ventilación y protecciones. En red se verifican vegetación, flechas, aisladores, conexiones, balance por fases, tierra y condición de postes. La periodicidad se fijará con fabricantes, criticidad y registros, en lugar de inventar un calendario obligatorio universal.

\section{Datos que cambian decisiones}
\begin{longtable}{L{3.3cm}L{5.8cm}L{4.8cm}}\toprule Comprobación & Trabajo necesario & Decisión que permite cerrar \\ \midrule
Caudal local & Aforos en estiaje y lluvias; curva de gasto, información SENAMHI/ANA autorizada & Capacidad y energía, caudal ecológico, necesidad de respaldo.\\
Topografía & Identificar toma y casa; nivelar lámina, eje, canal, puntos altos y río & Salto bruto/neto, longitudes, transitorios y emplazamiento.\\
Usuarios & Padrón de cuatro registros y medición de demanda & Cobertura real, calificación y número de grupos.\\
Central existente & Placas, inspección de turbina, generador, transformador y línea & Reutilización segura y alcance de rehabilitación.\\
Crecidas y sedimentos & Extremos, secciones, niveles, socavación y granulometría & Seguridad de bocatoma, obras y casa; desarenación.\\
Suelo y roca & EMS, taludes, ensayos y resistividad estratificada & Fundaciones, anclajes, sismo y tierra.\\
Fabricantes & Curvas, altitud, inercia, impedancias y tiempos & Regulación, protecciones, eficiencia y cierre.\\
Costos y permisos & Ofertas comparables, titularidad y consulta de autoridades & Inversión, financiamiento y ruta legal.\\ \bottomrule\end{longtable}
Las comprobaciones pendientes no impiden presentar el cálculo académico, pero impiden interpretar la propuesta como diseño listo para construcción. La incertidumbre puede cambiar el dimensionamiento; no se oculta tras decimales en un presupuesto o un mapa.

\section{Apoyo para exposición y archivos reproducibles}
Se entrega una exposición resumida con ubicación, evidencia, método, demanda, alternativas, hidráulica, regulación, seguridad, costos y conclusiones. El guion propone una intervención de aproximadamente 12 minutos y respuestas a preguntas previsibles. El estudiante debe explicar especialmente la diferencia entre nivel y caudal, datos simulados y medidos, demanda normativa y padrón, potencia nominal y potencia en estiaje.

Las entradas JSON/CSV, capítulos LaTeX, figuras y scripts conservan fuentes y huellas SHA-256. El constructor reproduce memoria, planos y exposición; R02 queda como antecedente histórico.

\chapter{Conclusiones y recomendaciones}
\section{Resultados del proyecto académico}
\begin{enumerate}
\item Se desarrolla un proyecto académico de diseño elaborado por estudiantes, con operación aislada. No se declara expediente constructivo ni permisos obtenidos.
\item Los antecedentes de 200 kW, 50 m y Banki siguen sin placa o medición independiente. Las casas de máquinas candidatas difieren 152 m. El salto bruto de gabinete \dato{Hb} m requiere nivelación.
\item PISCO 1981--2025 permite una transposición exploratoria por área geodésica, no un aforo oficial. La estabilidad geométrica 30/90 m no valida el recurso; el control de salida D8 falló. Se requiere calibración local.
\item Limbani y tres registros de Huancasayani suman \dato{PopBase} habitantes y \dato{VivsBase} viviendas ocupadas en 2017. Se proyectan \dato{Users} conexiones potenciales, sin sumar el distrito. Las bases rurales de 2023 producen demanda de \dato{Demand} kW.
\item Se compararon 150, 200, 400 y 600 kW. Para estudiar cobertura del escenario conservador se dimensionaron tres grupos hidráulicos de 200 kW, con caudal nominal \dato{Qnom} \(\mathrm{m^3/s}\), salto neto \dato{Hn} m y tubería nominal de 1000 mm por 350 m. Entregan aproximadamente \dato{Pusers} kW a usuarios en el punto nominal; su margen es reducido y requiere datos garantizados.
\item La potencia hidroeléctrica disponible no es permanente: el indicador diario de excedencia de 95\% del modelo es \dato{Pfirmstudy} kW en usuarios y el mínimo es nulo. Por ello no se garantiza suministro continuo de 150--200 kW en estiaje. La energía hidro servida es \dato{EHidro} MWh/año y el déficit \dato{Edef} MWh/año bajo el perfil supuesto. Se estudia respaldo modular y gestión de cargas.
\item El MOC exige coordinar ELC y cierre lento: a ocho segundos se estiman \dato{PNslow} bar. El primer pico brusco alcanza \dato{PNfast} bar y puede cavitar. PN40 es candidato, sin certificación contra vacío o colapso de cavidades.
\item El sistema eléctrico inicial es aislado, 380 V / 10 kV, con transformador principal de 800 kVA, seis subestaciones y auxiliares con neutro propio. La caída MT de estudio es \dato{DropMT}\%; la rama BT ejemplar tiene \dato{DropBT}\%. La corriente inicial máxima de barra estimada es \dato{IKmax} kA. La protección debe reconocer el decremento de falla de generadores; sobrecorriente simple y una resistencia de tierra aislada no bastan para demostrar seguridad.
\item La matriz separa obligaciones peruanas y recomendaciones. Se actualizan bases rurales 2023 y E.030 de 2026. Los 600 kW superan el umbral de concesión hidráulica de 500 kW; agua, ambiente y obras requieren procedimientos propios.
\item El presupuesto estimado de S/\dato{CapexM} millones presenta escenarios financieros desfavorables. Los 600 kW cubren mejor el pico, pero no demuestran inversión óptima. Se deben comparar rehabilitación y etapas con usuarios y recurso verificados.
\end{enumerate}

\section{Prioridades para continuar}
Antes de escoger capacidad se requieren aforos en captación, nivelación y padrón con demanda medida. Después corresponden crecidas, geotecnia, resistividad, curvas de fabricante, cotizaciones y permisos. Estos datos permitirán comparar rehabilitación, ampliación por etapas y otras fuentes. La eventual conexión al SEIN exige estudios y aprobación propios.

\appendix
\chapter{Matriz normativa y evidencia}
Los estados abiertos se refieren a verificaciones de una ejecución real o a datos no recuperados. Una comprobación académica con hipótesis no se registra como cumplimiento de campo.
\begingroup\footnotesize\renewcommand{\arraystretch}{1.08}\input{generado/matriz}\endgroup
\chapter{Fuentes, procedencia y alcance de lectura}
Consulta del 13--14 de octubre de 2026. N01--N28 identifican normas y referencias técnicas; H01--H18 identifican hidrología, censo, cartografía y antecedentes. Las fichas de normas pagadas se distinguen de textos completos. Los enlaces son consultables en esta edición digital; las copias y huellas originales se conservan en el proyecto.
\begingroup\small\input{generado/fuentes}\endgroup
\end{document}
